\documentclass[10pt,a4paper]{amsart}
\usepackage[margin=19mm]{geometry}
\usepackage{amsmath,amssymb,mathtools}
\usepackage[scr=rsfs,cal=euler]{mathalfa}
\usepackage{xfrac}
\usepackage[unicode,hidelinks,bookmarks=false]{hyperref}
\newcommand{\ie}{i.e.\ }
\newcommand{\cf}{cf.\ }
\newcommand{\resp}{respectively\ }
\newcommand{\Subsection}{\subsection}
\providecommand{\nobreakdash}{\nobreak\hbox{-}}
\setlength{\emergencystretch}{2em}
\multlinegap0pt
\usepackage{amssymb}
\newtheorem{lemma}{Lemma}
\newcommand{\Z}{\mathbb{Z}}
\title{Hyperbolicity of complements of orbits in Anosov flows}
\author{Sergio Fenley, Tali Pinsky, Mario Shannon}
\date{Source: arXiv:2607.28139v1 (CC BY 4.0)}
\begin{document}
\maketitle

\noindent\emph{Source and licence.} Hyperbolicity of complements of orbits in Anosov flows. Version arXiv:2607.28139v1 (CC BY 4.0), distributed by arXiv under the Creative Commons Attribution 4.0 International licence, http://creativecommons.org/licenses/by/4.0/, as recorded in the arXiv licence metadata for this exact version and shown on the abstract page https://arxiv.org/abs/2607.28139v1. Full licence notice: \texttt{licenses/arxiv-2607.28139-CC-BY-4.0.txt}.\par\smallskip

\noindent\textit{Editorial scope.} Counted unit: the article's lemma core (source0.tex lines 422--424). The article's complete original proof of it is delivered with it (source0.tex lines 427--441, one \texttt{proof} environment of 15 lines). No mathematics is written, repaired or re-derived: the statement and the proof are the article's own lines, in the article's own notation, and no retained line is substituted or re-flowed.  No further source-local context is needed: every cross-reference of every retained line resolves inside the two delivered spans.  Nothing was omitted from any retained span, so the package carries no omission record.  No retained line contains a citation, so the wrapper carries no bibliography and no bibliography entry.  No retained line contains a figure environment, an image-inclusion command or drawing code, so no image is delivered and the package declares no asset dependency.  Editorial formatting only, in addition: an AMS \texttt{amsart} preamble that restates the article's own declarations and macros used by the retained lines, namely the environments \texttt{lemma} on separate counters, together with the standard packages those lines need.  The retained environments print in this document as Lemma 1 in order of appearance, whereas the article prints the same items with the numbering of its own full document.  The complete native source of the article as distributed in the arXiv e-print for this exact version is bundled unchanged as \texttt{sources/206389/source0.tex}.
\begin{lemma}\label{lem:isotopic to core}
Let $\gamma$ be a periodic orbit of an Anosov flow $\phi$ on a closed 3-manifold $M$. If $\gamma$ is contained in a solid torus and is homotopic to a core $\alpha$ of the solid torus, then $\gamma$ and $\alpha$ are isotopic.
\end{lemma}

\begin{proof}
Let $V$ be the solid torus and $T$ its boundary. They will be fixed
throughout the proof. Consider the invariant manifold $L^s(\gamma) \cap V$ containing $\gamma$. Lifting to the universal cover $\tilde M$, $L^s(\gamma)$ lifts to a plane $\tilde L$, through which $\tilde\gamma$ is an infinite line, and it has at least two infinite components of the intersection with the infinite cylinder $\tilde T$, one on either side of $\gamma$.

 Isotope $\tilde L$ equivariantly in the complement of $\tilde\gamma$ so that it is transverse to $\tilde T$. Then, using that $M$ is irreducible, isotope it so that all
intersections with $T$ are not null homotopic. 
It follows that there is an annulus $S$, the quotient of the innermost component of $L_1\setminus \tilde T$ embedded in $V$ with core $\gamma$.

This annulus defines an isotopy from $\gamma$ to a curve $\gamma'\subset T$.  
Let $D$ be a meridian disk in $V$, well defined up to isotopy.
By the assumption that $\gamma$ and therefore $\gamma'$ is homotopic to the core of $V$, the arcs of $\gamma'\setminus D$ can be removed, innermost first, by an isotopy of $\gamma'$ within $T$, so that the resulting curve $\gamma''$ has a single intersection with $D$. By isotoping $\gamma''$ further, we may assume it is monotone with respect to the fibration of $T$ by curves parallel to $\partial D$.
Observe that, in an homology basis $\{l,m\}$ of $H_1(T)$ where $m=[\partial D]$ and $l$ is the class of a closed loop intersecting $\partial D$ only once, we have that $[\gamma"]=(1,n)$ for some $n\in\Z$. Hence, $\gamma''$ is isotopic to a core $\alpha$ of the solid torus $T$ by sliding it inwards the solid torus $V$ along meridians disks. Therefore, concatenating the all isotopies 
$$\gamma\to\gamma'\to\gamma ''\to\alpha$$ 
we conclude that $\gamma$ is isotopic to $\alpha$, as required.
\end{proof}

\end{document}
